%%%%%%%%%%%%%%%%%%%%%%%%%%%%%%%%%%%%%%%%%%%%%%%%%%%%%%%%%%%%%%%%%%%%%%
% Lớp 11 | Đạo hàm | Đạo hàm
% Dạng 1D7?1-4: Ý nghĩa Vật lý của đạo hàm
% Số câu: 31  (sinh tự động bởi tach_cau.py; từ nay thầy sửa trực tiếp file này)
%%%%%%%%%%%%%%%%%%%%%%%%%%%%%%%%%%%%%%%%%%%%%%%%%%%%%%%%%%%%%%%%%%%%%%

%%%%%%%%%%%%%%%% TRẮC NGHIỆM 4 PHƯƠNG ÁN %%%%%%%%%%%%%%%%

% nguồn: Ok_TracNghiemLop11_Lan1
\begin{ex}%[1D7N1-4]%[ma:730dcfa5]
	Cho chuyển động xác định bởi phương trình $S=t^3-3t^2-9t$, trong đó $t$ được tính bằng giây và $S$ được tính bằng mét. Tính vận tốc của chuyển động tại thời điểm gia tốc triệt  tiêu.
	 \choice
		{$3$  m/s}
		{\True  $-12$ m/s}
		{$12$  m/s}
		{$-3$  m/s}
	\loigiai{Gia tốc  của chuyển động là $a(t)=S''(t)=6t-6.$\\
		Gia tốc triệt tiêu khi $a(t)=0\Leftrightarrow t=1.$\\ 
		 Vận tốc tại thời điểm gia tốc triệt tiêu là $v(1)=S'(1)=-12$ m/s.
	} 
\end{ex}

% nguồn: Ok_Toan11_TracNghiem_Lan2
\begin{ex}%[1D7H1-4]%[ma:4c971407]
	Một chất điểm chuyển động  có phương trình chuyển động là $S(t)=\dfrac{1}{3}t^{3}-2t^{2}+3t-1$ (t được tính bằng giây, s tính bằng mét). Vận tốc tức thời của chất điểm tại thời điểm $t=4$ là
	\choice 
		 {$6$ (m/s)}
		{$4$ (m/s)} 
		{$5$  (m/s)}
		{\True $3$ (m/s)} 
	 \loigiai{
		\begin{itemize}
			 \item Ta có $v(t)=S'(t)=t^{2}-4t+3$.
			\item Vận tốc tức  thời của chất điểm tại thời điểm $t=4$ là $v(4) = 4^2-4(4)+3=3$ (m/s).
		\end{itemize}
	} 
\end{ex}

% nguồn: Ok_Toan11_TracNghiem_Lan2
\begin{ex}%[1D7H1-4]%[ma:c66b28f2]
	Cho chuyển động được xác định  bởi phương trình $s=2t^{3}+6t^{2}-t$, trong đó $t$ được tính bằng giây và $s$ được tính bằng mét. Vận tốc tức thời của chuyển động tại thời điểm $t=3$ s là
	\choice 
		{\True $89$  m/s}
		{$105$  m/s}
		{$48$ m/s} 
		{$20$  m/s}
	 \loigiai{
		\begin{itemize}
			 \item Ta có $v(t)=s'(t)=6t^{2}+12t-1$.
			\item Vận tốc tức  thời của chuyển động khi $t=3$ s là $v(3)=6 \cdot 3^{2}+12 \cdot 3-1=89$ (m/s).
		\end{itemize}
	 }
\end{ex}

% nguồn: Ok_TracNghiemLop11_Lan1
\begin{ex}%[1D7H1-4]%[ma:8c132cc2]
	Một chất điểm chuyển động có  phương trình $s=2t^2+3t$ ($t$ tính bằng giây, $s$ tính bằng mét). Vận tốc của chất điểm tại thời điểm $t_0=2$ (giây) bằng
	 \choice
		{$22$  m/s}
		 {$19$ m/s}
		{$9$ m/s} 
		{\True $11$  m/s}
	\loigiai{ 
		Phương trình vận tốc của chất điểm được xác định  bởi $v=s'=4t+3$.\\
		Suy ra vận  tốc của chất điểm tại thời điểm $t_0=2$ (giây) bằng $v(2)=4\cdot 2+3=11$.
	 }
\end{ex}

% nguồn: Ok_TracNghiemLop11_Lan1
\begin{ex}%[1D7H1-4]%[ma:0b303e70]
	Một chất điểm chuyển động có phương trình $s=2t^2+3t$ ($t$ tính bằng giây, $s$ tính bằng mét). Vận tốc của chất điểm tại  thời điểm $t_0=2$ (giây) bằng
	\choice 
		 {$22\left(m/s\right)$}
		{$19\left(m/s\right)$} 
		 {$9\left(m/s\right)$}
		{\True$11\left(m/s\right)$} 
	\loigiai{ 
		 Ta có $v(t)=s'(t)=4t+3$.\\
		Vận tốc của chất điểm tại thời điểm $t_0=2$  (giây) là $v\left(2\right)=11\left(m/s\right)$.
	 }
\end{ex}

% nguồn: Ok_Toan11_TracNghiem_Lan2
\begin{ex}%[1D7V1-4]%[ma:36054570]
	Một chất điểm chuyển động có quãng đường được cho bởi phương trình $s(t)=t^{3}-3t^{2}+5t+10$, trong đó $t>0$ với $t$ tính bằng giây và $s$ tính bằng mét.  Hỏi tại thời điểm vận tốc của vật đạt giá trị nhỏ nhất thì quãng đường vật đi được bằng bao nhiêu?
	 \choice
		{\True $13$ m} 
		{$3$  m}
		{$16$ m} 
		 {$10$ m}
	 \loigiai{
		\begin{itemize}
			\item Vận tốc  của chuyển động là $v(t)=s'(t)=3t^{2}-6t+5$.
			\item Ta có  $v(t)=3(t-1)^{2}+2 \ge 2$. Dấu \lq\lq$=$\rq\rq\, xảy ra khi $t=1$.
			\item Vận tốc  đạt giá trị nhỏ nhất tại $t=1$ s.
			\item Khi đó, quãng đường vật đi được là  $s(1)=1^3-3(1)^2+5(1)+10 = 13$ (m).
		\end{itemize}
	 }
\end{ex}

% nguồn: Ok_Toan11_TracNghiem_Lan2
\begin{ex}%[1D7V1-4]%[ma:43ba1185]
	Một chất điểm chuyển động theo quy luật  $S=6t^{2}-t^{3}$, vận tốc $v$ (m/s) của chuyển động đạt giá trị lớn nhất tại thời điểm $t$ (s) bằng
	 \choice
		{$4$ (s)} 
		{$12$  (s)}
		 {$6$ (s)}
		 {\True $2$ (s)}
	\loigiai{ 
		\begin{itemize}
			\item Ta có  $v(t)=S'(t)=12t-3t^{2}$.
			\item Đây là một parabol có bề lõm quay xuống, đạt giá trị lớn nhất tại đỉnh khi  $t=-\dfrac{12}{2(-3)}=2$ (s).
		\end{itemize}
	} 
\end{ex}

% nguồn: Ok_Toan11_TracNghiem_Lan2
\begin{ex}%[1D7V1-4]%[ma:69c88cf9]
	Một vật chuyển động theo quy luật $s(t)=4t^{2}-2t^{3}+5$ với $t$  là khoảng thời gian tính từ lúc vật bắt đầu chuyển động và $s$ là quãng đường vật đi được trong thời gian đó. Biết tại thời điểm $m$ thì vận tốc của chuyển động đạt giá trị lớn nhất là $n$ (m/s). Giá trị $T = mn$ bằng
	\choice 
		 {$\dfrac{4}{3}$}
		{\True  $\dfrac{16}{9}$}
		 {$\dfrac{2}{3}$}
		{$\dfrac{8}{9}$} 
	\loigiai{ 
		\begin{itemize}
			\item  Ta có $v(t)=s'(t)=8t-6t^2$.
			\item Vận tốc  của chuyển động đạt giá trị lớn nhất khi $t=-\dfrac{b}{2a}=-\dfrac{8}{2(-6)}=\dfrac{2}{3}$ (s). Vậy $m=\dfrac{2}{3}$.
			 \item $v_{\max}=v\left(\dfrac{2}{3}\right)=8\cdot\left(\dfrac{2}{3}\right)-6\cdot\left(\dfrac{2}{3}\right)^2=\dfrac{16}{3}-\dfrac{8}{3}=\dfrac{8}{3}$. Vậy $n=\dfrac{8}{3}$.
			\item  Vậy $T=mn=\dfrac{2}{3} \cdot \dfrac{8}{3}=\dfrac{16}{9}$.
		\end{itemize}
	} 
\end{ex}

% nguồn: Ok_Toan11_TracNghiem_Lan2
\begin{ex}%[1D7C1-4]%[ma:29dad855]
	\immini  {
	Một vật chuyển  động trong $1$ giờ với vận tốc $v$ phụ thuộc vào thời gian $t$ có đồ thị vận tốc như hình bên. Trong khoảng thời gian $1$ giờ kể từ khi bắt đầu chuyển động, đồ thị đó là một phần của đường parabol có đỉnh $I\left(\dfrac{1}{2};8\right)$ và trục đối xứng song song với trục tung. Tính gia tốc của vật lúc $t=0{,}25$ (h).
	 \choice
		{\True $16$  (km/h$^2$)}
		{$-16$  (km/h$^2$)}
		{$8$ (km/h$^2$)} 
		 {$-8$ (km/h$^2$)}
		} 
		 {
		\begin{tikzpicture}[transform shape]
				 \draw[-stealth] (-1,0)--(3,0) node[below]{$t$};
				\draw[-stealth] (0,-1)--(0,9) node [right]  {$v$};
				\draw[smooth] plot[domain=0:1]  (\x,{-32*(\x-0.5)^2+8});
				 \draw[dashed]
					(0,0) node[below left] {$O$} 
					(1,0) node[below]  {$1$}
					(.5,0)  node[below]{$\dfrac{1}{2}$}|-(0,8) node [left] {$8$};
			\end{tikzpicture}
		 }
	 \loigiai{
		\begin{itemize}
			\item Dựa vào đồ thị của hàm số, ta  suy ra $v(t)=a\cdot\left(x-\dfrac{1}{2}\right)^2+8$.
			 \item $v(1)=0\Leftrightarrow a\cdot\left(1-\dfrac{1}{2}\right)^2+8=0\Leftrightarrow a=-32$.
			\item Gia  tốc vật là $a(t)=v'(t)=-64\cdot\left(x-\dfrac{1}{2}\right)$.
			\item Lúc $t=0{,}25$ (h) thì gia tốc là $a(0{,}25)=-64\cdot\left(0{,}25-\dfrac{1}{2}\right)=16$  (km/h$^2$).
		\end{itemize}
	} 
\end{ex}


%%%%%%%%%%%%%%%% ĐÚNG/SAI %%%%%%%%%%%%%%%%

% nguồn: Ok_dungsai_lop11_Lan1
\begin{ex}%[1D7H1-4]%[ma:7cc1f010]
 Các khẳng định sau đây  đúng hay \textbf{sai}?
  \choiceTF
  {\True Nếu hàm số $s=f(t)$ biểu thị quãng đường di chuyển của vật theo thời gian $t$ thì $f'(t_0)$ biểu thị tốc  độ tức thời của chuyển động tại thời điểm $t_0$}
  {Hàm số $y=f(x)$ có đạo hàm tại $x_0$  là $f'(x_0)=\lim\limits_{x\to x_0}\dfrac{f(x)-f(x_0)}{x+x_0}$}
  {Hàm số $y=f(x)$ xác định trên khoảng $(a;b)$ có đại lượng $\Delta y=f(x_0)+f(x)$ gọi là số gia của biến tại $x_0$} 
  {\True Phương trình tiếp tuyến của đồ thị $(C): y=f(x)$ tại điểm $M(x_0;f(x_0))$  thuộc $(C)$ là $y-f(x_0)=f'(x_0)(x-x_0)$}
  \loigiai{
  \begin{enumerate}
   \item Nếu hàm số $s=f(t)$ biểu thị quãng đường di chuyển của vật theo thời gian $t$ thì $f'(t_0)$ biểu thị tốc độ tức  thời của chuyển động tại thời điểm $t_0$.
   \item Hàm số $y=f(x)$ có  đạo hàm tại $x_0$ là $f'(x_0)=\lim\limits_{x\to x_0}\dfrac{f(x)-f(x_0)}{x-x_0}$.
   \item Hàm số $y=f(x)$ xác định  trên khoảng $(a;b)$ có đại lượng $\Delta x=x-x_0$ gọi là số gia của biến tại $x_0$.
   \item Phương trình tiếp tuyến của đồ thị $(C): y=f(x)$ tại điểm $M(x_0;f(x_0))$ thuộc $(C)$  là $y-f(x_0)=f'(x_0)(x-x_0)$.
  \end{enumerate}
  }
\end{ex}

% nguồn: Ok_dungsai_lop11_Lan2
\begin{ex}%[1D7V1-4]%[ma:ecc856d2]
	Cho hàm số $y=f(x)=2x^3$ có đồ thị  $(C)$ và điểm $M$ thuộc $(C)$ có hoành độ $x_0=-1$.
	 \choiceTF
		{\True Hệ số góc của tiếp tuyến $\Delta$  của $(C)$ tại điểm $M$ bằng $6$}
		{\True Phương trình tiếp tuyến $\Delta$ đi qua  điểm $A(0;4)$}
		{Phương trình tiếp tuyến $\Delta$ cắt  đường thẳng $d\colon y=3x$ tại điểm có hoành độ bằng $4$}
		{\True Phương trình tiếp tuyến $\Delta$ vuông góc  với đường thẳng $\Delta'\colon y=-\dfrac{1}{6} x$}
	\loigiai{ 
		\begin{itemchoice}
			\itemch Hệ số góc của tiếp tuyến $\Delta$ của $(C)$ tại điểm $M$  là \[\lim\limits_{x\to -1}\dfrac{f(x)-f(-1)}{x+1}=\lim\limits_{x\to -1}\dfrac{2x^3+2}{x+1}=\lim\limits_{x\to -1}\left[2\left(x^2-x+1\right)\right]=6.\]
			Hệ số góc của tiếp  tuyến $\Delta$ tại điểm $M$ là $y'(-1)=6$.
			\itemch Tại điểm $x_0=-1$, ta có  $y_0=-2\Rightarrow M(-1;-2)$.\\
			Phương trình tiếp tuyến $\Delta$ của  $(C)$ tại $M$ là \[y-f(-1)=6(x+1)\Leftrightarrow y+2=6(x+1)\Leftrightarrow y=6x+4.\]
			Mà $4=6\cdot 0+4$ đúng nên tiếp tuyến $\Delta$  đi qua điểm $A(0;4)$.
			\itemch Phương trình hoành độ giao điểm của $\Delta$ và $d$  là $6x+4=3x\Leftrightarrow x=-\dfrac{4}{3}$.\\
			Vậy $\Delta$ cắt $d$ tại điểm $\left(-\dfrac{4}{3};-4\right)$. 
			\itemch Tiếp tuyến  $\Delta$ có hệ số góc $k_1=6$; đường thẳng $\Delta'$ có hệ số góc $k_2=-\dfrac{1}{6}$.\\
			Vì $k_1\cdot k_2=-1$ nên tiếp tuyến $\Delta$ vuông góc với  đường thẳng $\Delta'$.
		\end{itemchoice}
	} 
\end{ex}

% nguồn: Ok_dungsai_lop11_Lan2
\begin{ex}%[1D7V1-4]%[ma:b6cb65d9]
	Một vật chuyển động thẳng xác định bởi  phương trình $s=s(t)=t^3-3t^2-9t$ trong đó $t$ được tính bằng giây và $S$ được tính bằng mét.
	\choiceTF 
		{\True Vận tốc của vật tại thời điểm $t=5$\,s là  $36$\,m/s}
		{\True Gia tốc của vật tại thời điểm $t=3$\,s  là $6$\, m/s$^2$}
		 {\True Gia tốc tại thời điểm vận tốc triệt tiêu là $12$\,m/s$^2$}
		{Vận tốc tại thời điểm gia tốc triệt tiêu là $12$\,m/s} 
	 \loigiai{
		Tại $t_0 \in \mathbb{R}$ tuỳ ý, ta  có
		\begin{eqnarray*}
			s(t)-s\left(t_0 \right)&=&t^3-3t^2-9t-\left(t_0 ^3-3t_0 ^2-9t_0 \right)\\ 
			&=&\left(t-t_0 \right)\left(t^2+t\cdot t_0+t_0 ^2\right)-3\left(t-t_0 \right)\left(t+t_0 \right)-9\left(t-t_0  \right)\\
			&=&\left(t-t_0 \right)\left(t^2+t\cdot t_0+t_0  ^2-3t-3t_0-9\right).
		\end{eqnarray*}
		\begin{eqnarray*}
			\lim\limits_{t\to  t_0}\dfrac{s(t)-s\left(t_0 \right)}{t-t_0}&=&\lim\limits_{t\to t_0}\dfrac{\left(t-t_0 \right)\left(t^2+t\cdot t_0+t_0 ^2-3t-3t_0-9\right)}{t-t_0}\\
			&=&\lim\limits_{t\to t_0}\left(t^2+t\cdot t_0+t_0 ^2-3t-3t_0-9\right)\\ 
			 &=&3t_0 ^2-6t_0-9.
		\end{eqnarray*}
		Ta có $v(t)=s'(t)=3t^2-6t-9$.\\ 
		Tại $t_0 \in \mathbb{R}$ tùy ý, ta  có
		\begin{itemize}
			\item  $v(t)-v\left(t_0 \right)=3t^2-6t-9-\left(3t_0 ^2-6t_0-9\right)=\left(t-t_0 \right)\left(3t+3t_0-6\right)$.
			\item  $\lim\limits_{t\to t_0}\dfrac{v(t)-v\left(t_0 \right)}{t-t_0}=\lim\limits_{t\to t_0}\dfrac{\left(t-t_0 \right)\left(3t+3t_0-6\right)}{t-t_0}=\lim\limits_{t\to t_0}\left(3t+3t_0-6\right)=6t_0-6$.
		\end{itemize}
		Ta có $a(t)=v'(t)=6t-6$. 
		\begin{itemchoice}
			\itemch  Vận tốc của vật tại thời điểm $t=5$\,s là $v(5)=3\cdot 5^2-6\cdot 5-9=36$\,(m/s).
			\itemch Gia tốc của vật tại thời điểm $t=2$\,s là $a(2)=6\cdot 2-6=6$\,(m/s$^2$). 
			\itemch Khi vận tốc triệt tiêu, ta  có $3t^2-6t-9=0\Rightarrow t=3$ $(t>0)$.\\
			Khi đó gia tốc là $a(3)=6\cdot 3-6=12$\,m/s$^2$. 
			\itemch Khi gia  tốc triệt tiêu, ta có $6t-6=0\Rightarrow t=1$.\\
			Khi đó  vận tốc là $v(1)=3\cdot 1^2-6\cdot 1-9=-12$\,(m/s).
		\end{itemchoice}
	 }
\end{ex}

% nguồn: Ok_dungsai_lop11_Lan2
\begin{ex}%[1D7V1-4]%[ma:5dc88dbf]
	Cho hàm số $f(x)=\dfrac{2}{1-x}$ có đồ  thị $(C)$ và điểm $M(3;-1)\in (C)$.
	 \choiceTF
		{\True Hệ số góc  của tiếp tuyến $\Delta$ của $(C)$ tại điểm $M$ bằng $\dfrac{1}{2}$}
		{Phương  trình tiếp tuyến $\Delta$ song song với đường thẳng $y=-\dfrac{1}{2} x-\dfrac{5}{2}$}
		{\True Phương trình tiếp tuyến $\Delta$ vuông với đường  thẳng $y=-2x-\dfrac{5}{2}$}
		{\True Phương trình tiếp tuyến  $\Delta$ đi qua điểm $A\left(0;-\dfrac{5}{2} \right)$}
	 \loigiai{
		\begin{itemchoice}
			\itemch Hệ  số góc của tiếp tuyến $\Delta$ của $(C)$ tại điểm $M$ là \[\lim\limits_{x\to 3}\dfrac{f(x)-f(3)}{x-3}=\lim\limits_{x\to 3}\dfrac{\tfrac{2}{1-x}+1}{x-3}=\lim\limits_{x\to 3}\dfrac{-(x-3)}{(x-3)(1-x)}=\lim\limits_{x\to 3}\dfrac{-1}{1-x}=\dfrac{1}{2}.\]
			Hệ  số góc của tiếp tuyến $\Delta$ của $(C)$ tại điểm $M$ là $y'(3)=\dfrac{1}{2}$.
			\itemch Ta  có
			\begin{itemize}
				\item Hệ  số góc của tiếp tuyến của $(C)$ tại điểm $M$ là $k_1=y'(3)=\dfrac{1}{2}$.
				\item Hệ số góc của đường  thẳng $d\colon y=-\dfrac{1}{2} x-\dfrac{5}{2}$ là $k_2=-\dfrac{1}{2}$.
			\end{itemize}
			Vì $k_1\neq k_2$ nên tiếp tuyến $\Delta$ không song  song với đường thẳng $d$.
			\itemch Ta  có
			\begin{itemize}
				\item  Hệ số góc của tiếp tuyến $\Delta$ của $(C)$ tại điểm $M$ là $k_1=y'(3)=\dfrac{1}{2}$.
				\item  Hệ số góc của đường thẳng $\Delta'\colon y=-2x-\dfrac{5}{2}$ là $k_3=-2$.
			\end{itemize}
			Vì $k_1\cdot k_3=-1$ nên tiếp tuyến của  $(C)$ tại $M$ song song với đường thẳng $\Delta'$.
			\itemch Phương trình tiếp tuyến $\Delta$ của $(C)$ tại $M(3;-1)$ là \[y-(-1)=\dfrac{1}{2}  (x-3)\Leftrightarrow y=\dfrac{1}{2} x-\dfrac{5}{2}.\]
			Vì $-\dfrac{5}{2}=\dfrac{1}{2}\cdot 0-\dfrac{5}{2}$ nên tiếp tuyến $\Delta$ đi  qua điểm $A\left(0;-\dfrac{5}{2} \right)$.
		\end{itemchoice}
	} 
\end{ex}

% nguồn: Ok_dungsai_lop11_Lan2
\begin{ex}%[1D7V1-4]%[ma:7dc80a7e]
	Cho hàm số $y=x^2+3x+1$ có đồ  thị $(C)$. Viết được phương trình tiếp tuyến $\Delta$ của $(C)$ tại giao điểm $M$ của $(C)$ với trục tung.
	\choiceTF 
		{\True Hệ số góc của tiếp tuyến $\Delta$ tại  điểm $M$ bằng $3$}
		{Phương trình tiếp tuyến $\Delta$  đi qua điểm $A(1;3)$}
		{\True Phương trình tiếp tuyến $\Delta$ cắt đường  thẳng $y=2x+1$ tại điểm có hoành độ bằng $0$}
		{\True Phương trình tiếp tuyến $\Delta$ vuông góc với đường thẳng $y=-\dfrac{1}{3} x+1$} 
	\loigiai{ 
		\begin{itemchoice}
			 \itemch
			\begin{itemize}
				\item Tại giao điểm $M$ của $(C)$ với trục tung, ta có $x_0=0\Rightarrow y_0=1\Rightarrow M(0;1)$. 
				\item Hệ số góc tiếp tuyến  $\Delta$ tại điểm $M$ là \[f'(0)=\lim\limits_{x\to 0}\dfrac{f(x)-f(0)}{x-0}=\lim\limits_{x\to 0}\dfrac{x^2+3x}{x}=\lim\limits_{x\to 0}(x+3)=3.\]
			\end{itemize}
			Hệ số góc của tiếp tuyến $\Delta$ của $(C)$  tại điểm $M$ là $y'(0)=3$.
			\itemch Phương trình tiếp tuyến $\Delta$ của $(C)$ tại điểm $M$ là  $y=3x+1$.\\
			Vì $3=3\cdot 0+1$ sai nên tiếp tuyến của $(C)$ tại  điểm $M$ không đi qua điểm $A$.
			\itemch Phương trình hoành độ giao  điểm của $\Delta$ và đường thẳng $y=2x+1$ là \[3x+1=2x+1\Leftrightarrow x=0.\]
			Suy ra tiếp tuyến $\Delta$ cắt  đường thẳng $y=2x+1$ tại điểm có hoành độ bằng $0$.
			\itemch \begin{itemize}
			\item Hệ số góc của tiếp tuyến $\Delta$ của $(C)$ tại  điểm $M$ là $y'(0)=3$.
			\item Hệ số góc của  đường thẳng $y=-\dfrac{1}{3} x+1$ là $k=-\dfrac{1}{3}$.
		\end{itemize}
		Vì $y'(0)\cdot k=-1$ nên tiếp tuyến $\Delta$ vuông  góc với đường thẳng $y=-\dfrac{1}{3} x+1$.
	\end{itemchoice}
 }
\end{ex}

% nguồn: Ok_dungsai_lop11_Lan2
\begin{ex}%[1D7V1-4]%[ma:a4479997]
	Viết được phương trình tiếp tuyến $\Delta$ của đồ thị hàm  số $y=\dfrac{4}{x-1}$ tại điểm có hoành độ $x_0=-1$.
	 \choiceTF
		{Hệ số góc tiếp tuyến $\Delta$  bằng $1$}
		{Phương trình tiếp tuyến  $\Delta$ đi qua điểm $M(-1;2)$}
		{\True Phương trình tiếp tuyến $\Delta$ cắt đường thẳng $y=2x+1$ tại  điểm có hoành độ bằng $\dfrac{4}{3}$}
		{\True Phương trình tiếp  tuyến $\Delta$ vuông góc với đường thẳng $y=x+1$}
	 \loigiai{
		\begin{itemchoice}
			\itemch Hệ  số góc tiếp tuyến $\Delta$ là
			\[f'(-1)=\lim\limits_{x\to -1}\dfrac{f(x)-f(-1)}{x+1}=\lim\limits_{x\to  -1}\dfrac{\tfrac{4}{x-1}+2}{x+1}=\lim\limits_{x\to -1}\dfrac{2(x+1)}{(x+1)(x-1)}=\lim\limits_{x\to -1}\dfrac{2}{x-1}=-1.\]
			Hệ số góc tiếp tuyến $\Delta$ là $y'(0)=-1$. 
			\itemch Tại điểm  $x_0=-1$, ta có $y_0=-2$.\\
			Khi đó,  phương trình tiếp tuyến $\Delta$ là $y=-x-3$.\\
			Vì $2=-(-1)-3$ sai nên tiếp tuyến  $\Delta$ không đi qua điểm $M(-1;2)$.
			\itemch Phương trình hoành độ  giao điểm của $\Delta$ và đường thẳng $y=2x+1$ là \[-x-3=2x+1\Leftrightarrow x=-\dfrac{4}{3}.\]
			 Suy ra $\Delta$ cắt đường thẳng $y=2x+1$ tại điểm có hoành độ bằng $-\dfrac{4}{3}$.
			\itemch \begin{itemize}
			\item Hệ số góc của  tiếp tuyến $\Delta$ là $y'(0)=-1$.
			\item Hệ  số góc của đường thẳng $y=x+1$ là $k=1$.
		\end{itemize}
		Vì $y'(0)\cdot k=-1$ nên tiếp tuyến  $\Delta$ vuông góc với đường thẳng $y=x+1$.
	\end{itemchoice}
 }
\end{ex}


%%%%%%%%%%%%%%%% TRẢ LỜI NGẮN %%%%%%%%%%%%%%%%

% nguồn: Ok_TraLoiNgan_Lop11_Lan1
\begin{ex}%[1D7H1-4]%[ma:42972787]
	Một quả bóng được thả rơi tự do từ đài quan sát trên sân thượng của toà nhà Landmark 81 (Thành phố Hồ Chí Minh) cao $461{,}3 m$ xuống mặt đất, với phương trình chuyển động $s(t)=4{,}9 t^2$. Tính  vận tốc của quả bóng khi nó chạm đất, bỏ qua sức cản không khí. (Đơn vị m/s, kết quả gần đúng làm tròn đến hàng phần chục)
	\shortans{95,1} 
	 \loigiai{
		Với bất kì $t_0$,  ta có
		 $$
		 s'\left(t_0\right)
		=\lim\limits_{t \rightarrow t_0}  \dfrac{s(t)-s\left(t_0\right)}{t-t_0}
		=\lim\limits_{t \rightarrow  t_0} \dfrac{4{,}9 t^2-4{,}9 t_0^2}{t-t_0}
		=\lim\limits_{t \rightarrow t_0} 4{,}9\left(t+t_0\right)=9{,}8  t_0 .
		$$ 
		Do đó, vận tốc của quả  bóng tại thời điểm $t$ là $v(t)=s'(t)=9{,}8 t$.
		 \\
		Mặt khác, vì chiều cao của toà tháp là $461{,}3$ m  nên quả bóng sẽ chạm đất tại thời điểm $t_1$. Từ đó, ta có
		 $$
		4{,}9  t_1^2=461{,}3
		 \quad\text{hay}\quad
		 t_1=\sqrt{\dfrac{461{,}3}{4{,}9}} \, \text{giây}.
		 $$
		Vậy vận tốc của quả bóng khi  nó chạm đất là
		 $$
		 v\left(t_1\right)
		 =9{,}8 t_1
		=9{,}8  \cdot \sqrt{\dfrac{461{,}3}{4{,}9}}
		\approx  95{,}1\, \text {m/s}
		$$ 
	 }
\end{ex}

% nguồn: Ok_TraLoiNgan_Lop11_Lan1
\begin{ex}%[1D7H1-4]%[ma:22c229a6]
	Một vật được phóng theo phương thẳng đứng lên trên từ mặt đất, biết độ cao $h$ của nó (tính bằng  mét) sau $t$ giây được cho bởi phương trình $h(t)=24{,}5 t-4{,}9 t^2$. Tìm vận tốc của vật khi nó chạm đất. (đơn vị m/s, làm tròn đến hàng phần mười)
	 \shortans{24,5}
	\loigiai{ 
		 Khi vật chạm đất thì
		 $$
		 h=0
		\Leftrightarrow 24{,}5 t-4{,}9 t^2=0 
		\Rightarrow  t=5.
		$$ 
		Ta có $v(t)=h'(t)=24{,}5-9{,}8 t$ nên tốc độ của  vật tại thời điểm nó chạm đất $t=5$ là
		 $$
		 v(5)=\left|h'(5)\right|=|24{,}5-9{,}8.5|=24{,}5 \, \text{m/s}.
		 $$
	 }
\end{ex}

% nguồn: Ok_TraLoiNgan_Lop11_Lan2
\begin{ex}%[1D7H1-4]%[ma:88d690da]
	Một chất điểm chuyển động thẳng xác định bởi phương trình $s(t)=\dfrac{1}{2} t^2$,  trong đó $t$ là thời gian tính bằng giây và $s$ là quãng đường đi được trong $t$ giây tính bằng mét. Tính vận tốc tức thời của chất điểm tại $t=5$\,s.
	 \shortans{5}
	 \loigiai{
		Vận tốc tức thời của chất điểm tại $t=5$\,s  là $v(t)=s'(t)=t$.\\
		 Vậy $v(5)=s' (5)=5$\,(m/s).
	} 
\end{ex}

% nguồn: Ok_TraLoiNgan_Lop11_Lan1
\begin{ex}%[1D7V1-4]%[ma:b0e21049]
	Cho biết điện lượng truyền trong dây dẫn theo thời gian biểu thị bởi hàm số $Q(t)=2  t^2+t$, trong đó $t$ được tính bằng giây và $Q$ được tính theo coulomb. Cường độ dòng điện tại thời điểm $t=4(s)$ là bao nhiêu C/s?
	 \shortans{17}
	\loigiai{ 
		Ta có 
		$Q'(t)=\left(2 t^2+t\right)'=4 t+1$ nên cường độ dòng điện tại thời điểm  $t=4(s)$ là
		$$ 
		 Q'(4)=4 \cdot 4+1=17\,\text{C/s}.
		$$ 
	} 
\end{ex}

% nguồn: Ok_TraLoiNgan_Lop11_Lan1
\begin{ex}%[1D7V1-4]%[ma:e5de6269]
	Một chất điểm  chuyển động thẳng xác định bởi phương trình $s(t)=\dfrac{1}{2} t^2$, trong đó $t$ là thời gian tính bằng giây và $s$ là quãng đường đi được trong $t$ giây tính bằng mét. Vận tốc tức thời của chất điểm tại $t=5$ là bao nhiêu m/giây.
	 \shortans{5}
	\loigiai{ 
		Vận tốc tức  thời của chất điểm tại $t=5$ là
		$$ 
		s'(5)=\lim\limits_{t \rightarrow  5} \dfrac{s(t)-s(5)}{t-5}
		 =\lim\limits_{t \rightarrow 5} \dfrac{\dfrac{1}{2} t^2-\dfrac{25}{2}}{t-5}
		=\lim\limits_{t  \rightarrow 5} \dfrac{\dfrac{1}{2}(t-5)(t+5)}{t-5}
		 =\lim\limits_{t \rightarrow 5} \dfrac{1}{2}(t+5)
		 =5.
		$$ 
		Vậy $v(5)=s'(5)=5$  m/s.
	} 
\end{ex}

% nguồn: Ok_TraLoiNgan_Lop11_Lan1
\begin{ex}%[1D7V1-4]%[ma:2c77200f]
	Một viên đạn được bắn lên trời từ một vị trí cách mặt đất $1000$m theo phương thẳng đứng với vận tốc ban đầu $v_0=294$m/s (bỏ qua sức cản của không khí).  Hỏi khi viên đạn đạt độ cao lớn nhất và sẽ bắt đầu rơi thì viên đạn cách mặt đất bao nhiêu mét?
	\shortans{5410} 
	 \loigiai{
	 \immini{
		Chọn trục $Oy$ theo phương thẳng đứng, chiều dương hướng từ mặt đất lên trời, gốc $O$ ở mặt đất và cách $A$ là vị trí  viên đạn được bắn lên, gốc thời gian (tức lúc $t=0$) được tính từ vị trí $A$ như hình bên; khi đó chuyển động của viên đạn là chuyển động biến đổi đều với vận tốc ban đầu $v_0=294$ m/s và với gia tốc. (Gia tốc nhận giá trị âm vì véc-tơ gia tốc ngược chiều dương của trục $Oy$. Phương trình chuyển động của viên đạn là
		$y=1000+294t-4{,}9t^2$ \\ 
		Ta có $v(t)=y'=294-9{,}8t$. \\ 
		Viên đạn đạt độ cao lớn nhất và sẽ bắt đầu rơi \\$v(t)=0 \Leftrightarrow  294-9{,}8t=0 \Leftrightarrow t=30$ (s). \\
		Khi đó viên đạn  cách mặt đất là\\ $y(30)=1000+294.30-4{,}{9.30}^2=5410$ (m).}{
		\begin{tikzpicture}[line cap=round,line join=round,scale=1,>=stealth]
				\coordinate[label=below:{$O$}]  (O) at (0,0);
				\coordinate[label=right:{$A$}] (A) at (0,1); 
				 \coordinate[label=right:{$B$}] (B) at (0,5.41);
				\coordinate  (C) at (-2,0);
				\coordinate (D) at  (2,0);
				\draw [->]  (0,0)--(0,6) node[right]{$y$};
				 \draw (C)--(D);
				\foreach \x  in {-2,-1.5,-1,-0.5,0,0.5,1,1.5,2}{
					\draw[-]  (\x-.1,-.2)--(\x,0);}
				\draw (0,1) node [left]{$1000$}; 
				\draw  (0,5.41) node [left]{$5410$};
				\draw  (0.5,5.41) node [right]{$(t=30)$};
				\draw (0.5,1)  node [right]{$(t=0)$};
				\foreach \p  in {O,A,B}
				\fill  (\p) circle (1pt);
			\end{tikzpicture}
		 }
	 }
\end{ex}

% nguồn: Ok_TraLoiNgan_Lop11_Lan2
\begin{ex}%[1D7V1-4]%[ma:e18ca3b1]
	Một chất điểm chuyển động có phương trình chuyển động là $s(t)=t^2+4t+6$ ($t$ được tính bằng giây, $s$ được tính bằng mét). Tính vận tốc tức thời của chuyển  động tại thời điểm $t=5$\,s.
	 \shortans{14}
	 \loigiai{
		Vận tốc tức thời của chuyển động tại thời điểm $t=5$ là 
		\begin{eqnarray*}
			v(5)=s'(5)&=&\lim\limits_{t\to 5} \dfrac{s(t)-s(5)}{t-5}\\ 
			 &=&\lim\limits_{t\to 5}\dfrac{t^2+4t+6-51}{t-5}\\
			&=&\lim\limits_{t\to  5} \dfrac{t^2+4t-45}{t-5}\\
			 &=&\lim\limits_{t\to 5}\left(t+9\right)\\
			&=&14\text{  (m/s)}.
		\end{eqnarray*}
	} 
\end{ex}

% nguồn: Ok_TraLoiNgan_Lop11_Lan2
\begin{ex}%[1D7V1-4]%[ma:95b661f9]
	Một chất điểm chuyển động thẳng xác định bởi phương trình $s(t)=\dfrac{1}{2} t^2$, trong đó $t$ là thời gian tính bằng giây và $s$ là quãng đường đi được trong  $t$ giây tính bằng mét. Vận tốc tức thời của chất điểm tại $t=5$\,s đạt bao nhiêu m/s?
	 \shortans{5}
	 \loigiai{
		Vận tốc tức thời của chất  điểm tại $t=5$ là
		\begin{eqnarray*}
			 s'(5)&=&\lim\limits_{t\to 5}\dfrac{s(t)-s(5)}{t-5}\\
			&=&\lim\limits_{t\to 5}\dfrac{\dfrac{1}{2}  t^2-\dfrac{25}{2}}{t-5}\\
			 &=&\lim\limits_{t\to 5}\dfrac{\dfrac{1}{2} \left(t-5\right)\left(t+5\right)}{t-5}\\
			&=&\lim\limits_{t\to 5}\dfrac{1}{2} \left(t+5\right)\\ 
			&=&5. 
		\end{eqnarray*}
		 Vậy $v(5)=s'(5)=5$ (m/s).
	 }
\end{ex}

% nguồn: Ok_TraLoiNgan_Lop11_Lan2
\begin{ex}%[1D7V1-4]%[ma:d4f7cdea]
	Một quả bóng được thả rơi tự do từ đài quan sát trên sân thượng của toà nhà Landmark $81$ (Thành phố Hồ Chí Minh) cao $461{,}3$\,m xuống mặt đất, với phương trình chuyển động $s(t)=4{,}9t^2$. Tính vận tốc của quả bóng khi nó chạm đất, bỏ qua sức cản không khí. (Đơn vị m/s, kết quả  gần đúng làm tròn đến hàng phần chục).
	 \shortans{95,1}
	\loigiai{ 
		Với  bất kì $t_0$, ta có
		\[s' \left(t_0 \right)=\lim\limits_{t\to t_0} \dfrac{s(t)-s\left(t_0  \right)}{t-t_0}=\lim\limits_{t\to t_0}\dfrac{4{,}9t^2-4{,}9t_0^2}{t-t_0}=\lim\limits_{t\to t_0} 4{,}9\left(t+t_0 \right)=9{,}8t_0.\]
		Do đó, vận tốc của quả bóng tại thời điểm $t$ là  $v(t)=s' (t)=9{,}8t$.\\
		Mặt khác, vì chiều cao của toà tháp là $461{,}3\; m$ nên quả bóng sẽ chạm  đất tại thời điểm $t_1$.\\
		Từ đó,  ta có $4{,}9t_1^2=461{,}3\Leftrightarrow t_1=\sqrt{\dfrac{461{,}3}{4{,}9}}$ (giây).\\
		Vậy vận tốc của quả bóng  khi nó chạm đất là
		\[v\left(t_1 \right)=9{,}8t_1=9{,}8\cdot \sqrt{\dfrac{461{,}3}{4{,}9}} \approx  95{,}1\text{ (m/s)}.
		 \]
	} 
\end{ex}


%%%%%%%%%%%%%%%% TỰ LUẬN %%%%%%%%%%%%%%%%

% nguồn: Ok_TuanLuan_Toan11_Lan1
\begin{bt}%[1D7H1-4]%[ma:68062199]
	Một chất điểm chuyển động theo phương trình $s(t)=3 \sin 2 t+2 \cos 2 t$, trong đó $t$ là thời gian tính bằng giây và $s(t)$ là quãng đường chuyển động được của chất điểm trong $t$ giây tính bằng mét. Tính gia tốc của chất điểm đó khi  $t=\dfrac{\pi}{4}$.
	 \loigiai{ \\
		Ta có $s'(t)=3(\sin 2 t)'+2(\cos 2  t)'=6 \cos 2 t-4 \sin 2 t$.\\
		Và $s''(t)=6(\cos 2 t)'-4(\sin 2 t)'=-12 \sin 2 t-8 \cos  2 t$.\\
		Gia tốc  của chất điểm tại thời điểm $t=\dfrac{\pi}{4}$ là \\
		$a\left(\dfrac{\pi}{4}\right)=s'' \left(\dfrac{\pi}{4}\right)=-12\left[\sin  \left(2 \cdot \dfrac{\pi}{4}\right)\right]-8\left[\cos \left(2 \cdot \dfrac{\pi}{4}\right)\right]=-12$.
	} 
\end{bt}

% nguồn: Ok_TuanLuan_Toan11_Lan2
\begin{ex}%[1D7H1-4]%[ma:efdcdbae]
	Một chất điểm chuyển động  theo phương trình $s(t)=3 \sin 2 t+2 \cos 2 t$, trong đó $t$ là thời gian tính bằng giây và $s(t)$ là quãng đường chuyển động được của chất điểm trong $t$ giây tính bằng mét. Tính gia tốc của chất điểm đó khi $t=\dfrac{\pi}{4}$.
	\loigiai{ \\ 
		Ta có $s'(t)=3(\sin 2 t)'+2(\cos 2 t)'=6 \cos 2  t-4 \sin 2 t$.\\
		Và  $s''(t)=6(\cos 2 t)'-4(\sin 2 t)'=-12 \sin 2 t-8 \cos 2 t$.\\
		Gia tốc của chất điểm tại thời điểm $t=\dfrac{\pi}{4}$ là \\ 
		$a\left(\dfrac{\pi}{4}\right)=s'' \left(\dfrac{\pi}{4}\right)=-12\left[\sin \left(2 \cdot \dfrac{\pi}{4}\right)\right]-8\left[\cos  \left(2 \cdot \dfrac{\pi}{4}\right)\right]=-12$.
	 }
\end{ex}

% nguồn: Ok_TuanLuan_Toan11_Lan2
\begin{ex}%[1D7H1-4]%[ma:ff745530]
	Một chất điểm chuyển động thẳng biến đổi  đều với phương trình $ s=2 t^{2}+t-1$ m.
	\begin{enumerate}
		\item Tìm vận tốc tức thời của vật tại thời điểm $ t=2$  giây.
		\item Tìm vận tốc trung  bình của chất điểm trong khoảng thời gian từ $ t=0 $ tới $ t=2$ giây.
	\end{enumerate}
	 \loigiai{
		\begin{enumerate}
			\item Ta có  $ v=s'=4 t+1 $.\\
			Vận tốc tức thời của vật tại thời điểm $ t=2 $ s là $ 4\cdot  2+1=9$ m/s.
			\item Trong khoảng thời gian từ $ t=0 $ tới $ t=2$ s thì chất điểm di chuyển được quãng đường: 
			\begin{eqnarray*}
				s(2) - s(0)  &=& \left(2(2)^2 + 2 -1\right) - \left(2(0)^2 + 0 - 1\right) \\
				 &=& 9 - (-1) = 10.
			\end{eqnarray*}
			Suy ra vận tốc trung bình của chất  điểm trong khoảng thời gian $ 2$ s kể từ thời điểm $ t=0 $ là
			 \[
			 \overline{v}=\dfrac{\Delta s}{\Delta t}=\dfrac{10}{2}=5.
			\] 
			Vậy vận  tốc cần tìm là $5$ m/s.
		\end{enumerate}
	 }
\end{ex}

% nguồn: Ok_TuanLuan_Toan11_Lan1
\begin{bt}%[1D7V1-4]%[ma:75650c6d]
	Một chất điểm chuyển động thẳng xác định bởi phương trình $s(t)=\dfrac{1}{2} t^2$, trong đó $t$ là thời gian tính  bằng giây và $s$ là quãng đường đi được trong $t$ giây tính bằng mét. Vận tốc tức thời của chất điểm tại $t=5$ bằng bao nhiêu?
	\loigiai{ 
		Vận tốc tức thời của chất điểm  tại $t=5$ là
		\[s'(5) = \displaystyle  \lim \limits_{t \to 5} \dfrac{s(t)-s(5)}{t-5} = \lim \limits_{t \to 5} \dfrac{\tfrac{t^2}{2} - \tfrac{25}{2}}{t-5} =\lim \limits_{t \to 5} \dfrac{t+5}{2} =5.\]
		Vậy  $v(5)=s'(5)=5(\text{m/s})$.
	} 
\end{bt}

% nguồn: Ok_TuanLuan_Toan11_Lan1
\begin{bt}%[1D7V1-4]%[ma:5c61091b]
	Cân nặng trung bình của một em bé trong độ tuổi từ $0$ đến $36$ tháng có thể được tính gần đúng bởi hàm số $w(t) = 0{,}00076t^3 - 0{,}06t^2 + 1{,}8t + 8{,}2$. Trong đó $t$ được tính bằng tháng và $w$ được tính bằng  pound. Tính tốc độ thay đổi cân nặng của em bé đó tại thời điểm $15$ tháng tuổi.
	\loigiai{ 
		$w'(t) =  0{,}0028t^2 - 0{,}12t + 1{,}8$.
		\\Tốc độ thay đổi cân nặng của em  bé tại thời điểm $15$ tháng tuổi
		\\$w'(15) = 0{,}0028\cdot 15^2 -  0{,}12\cdot 15 + 1{,}8 = 0{,}513$.
	 }
\end{bt}

% nguồn: Ok_TuanLuan_Toan11_Lan2
\begin{ex}%[1D7V1-4]%[ma:003d956b]
	Một chất điểm chuyển động thẳng xác định bởi phương trình $s(t)=\dfrac{1}{2} t^2$, trong đó $t$ là thời gian tính bằng giây và $s$  là quãng đường đi được trong $t$ giây tính bằng mét. Vận tốc tức thời của chất điểm tại $t=5$ bằng bao nhiêu?
	\loigiai{ 
		Vận tốc tức thời của chất điểm  tại $t=5$ là
		\[s'(5) = \displaystyle \lim \limits_{t \to 5}  \dfrac{s(t)-s(5)}{t-5} = \lim \limits_{t \to 5} \dfrac{\tfrac{t^2}{2} - \tfrac{25}{2}}{t-5} =\lim \limits_{t \to 5} \dfrac{t+5}{2} =5.\]
		Vậy $v(5)=s'(5)=5(\text{m/s})$. 
	} 
\end{ex}

% nguồn: Ok_TuanLuan_Toan11_Lan2
\begin{ex}%[1D7V1-4]%[ma:37e0fcb4]
	Cân nặng trung bình của một em bé trong độ tuổi từ $0$ đến $36$ tháng có thể được tính gần đúng bởi hàm số $w(t) = 0{,}00076t^3 - 0{,}06t^2 + 1{,}8t + 8{,}2$. Trong đó $t$ được tính bằng tháng và $w$ được tính bằng pound. Tính tốc  độ thay đổi cân nặng của em bé đó tại thời điểm $15$ tháng tuổi.
	\loigiai{ 
		$w'(t) = 0{,}0028t^2 - 0{,}12t  + 1{,}8$.
		\\Tốc độ thay đổi cân nặng của em bé tại thời điểm $15$  tháng tuổi
		\\$w'(15) = 0{,}0028\cdot 15^2 - 0{,}12\cdot  15 + 1{,}8 = 0{,}513$.
	} 
\end{ex}
